\unit{Exposé 190 -- Descent Technique I}
\sid{190}
\src{299}
\seminaren{12th year, 1959/60, no. 190\\December 1959}
\paperhead{DESCENT TECHNIQUE AND EXISTENCE THEOREMS IN ALGEBRAIC GEOMETRY\\
I. GENERALITIES. DESCENT BY FAITHFULLY FLAT MORPHISMS}%
  {by Alexander GROTHENDIECK}

From a technical point of view, the present exposé and those that follow it may
be regarded as variations on Hilbert's celebrated “Theorem 90.” The
introduction of the descent method into algebraic geometry seems to be due to
A. Weil, under the name “descent of the base field.” Weil restricts himself to
finite separable field extensions. The case of radicial extensions of height 1
was subsequently treated by P. Cartier. For lack of the language of schemes,
and more particularly for lack of nilpotent elements in the rings one allowed
oneself to consider, Cartier had been unable to formulate the essential
identity between these two cases.

At present, it seems that the general descent technique to be presented
(combined, where appropriate, with the fundamental theorems of “formal
geometry,” cf. [3]) underlies most existence theorems in algebraic
geometry.\footnote{The 1962 erratum calls this formulation excessive: it is
true to a large extent for the nonprojective techniques treated in the first
two exposés of the TDTE I--VI series, but not for the projective techniques
(TDTE IV, V, and VI).} It should moreover be noted that this technique can
certainly be carried over to “analytic geometry,” and one may hope that in the
near future specialists will be able to prove the “analytic” analogues of the
existence theorems in formal geometry that will be given in Exposé II. In any
case, the recent work of Kodaira--Spencer, whose methods seem ill-suited to the
definition and study of “moduli varieties” near their singular points, shows
quite clearly the need for methods closer to scheme theory (which must of
course be supplemented by transcendental techniques).

In the present Exposé I we shall treat the most elementary case of descent
(indicated in the title). The applications of Theorems 1, 2, and 3 below are
nevertheless already very numerous. We have confined ourselves to giving a few
of them as examples, without attempting to state them here in the greatest
possible generality.

We shall freely use the language of schemes, for which we refer to the exposé
already cited, as well as to [2]. We note explicitly, however, that the
preschemes considered in the present exposé are not necessarily Noetherian, and
the

\src{300}
morphisms are not necessarily of finite type. Thus, if $A$ is a Noetherian
local ring with completion $\widehat A$, we shall have to consider the
non-Noetherian ring $\widehat A\otimes_A\widehat A$ and the morphisms of affine
schemes corresponding to the inclusions among the rings under consideration.

\fsection{A}{Preliminaries on categories}

\subsection*{1. Fibered categories, descent data, and
$\mathcal F$-descent morphisms}
\addcontentsline{toc}{subsection}{1. Fibered categories and descent data}

\result{a. DEFINITION}{1.1} A \emph{fibered category} $\mathcal F$ over a base
$\mathcal C$ consists of a category $\mathcal C$, a category $\mathcal F_X$
for every $X\in\mathcal C$, a functor
$f^*:\mathcal F_Y\longrightarrow\mathcal F_X$ for every
$\mathcal C$-morphism $f:X\longrightarrow Y$, also denoted by
\[
  f^*(\xi)=\xi\times_Y X
\]
for $\xi\in\mathcal F_Y$ (with $X$ regarded as an “object of $\mathcal C$
over $Y$,” i.e. equipped with the morphism $f$), and finally, for two
composable morphisms $X\xrightarrow{f}Y\xrightarrow{g}Z$, an isomorphism of
functors
\[
  c_{f,g}:(gf)^*\longrightarrow f^*g^*,
\]
these data being subject to the following conditions:

\begin{enumerate}
\item[(i)] $\operatorname{id}_X^*=\operatorname{id}$;
\item[(ii)] $c_{f,g}$ is the identity isomorphism if either $f$ or $g$ is an
identity isomorphism;
\item[(iii)] given three composable morphisms
$X\xrightarrow{f}Y\xrightarrow{g}Z\xrightarrow{h}T$, the following diagram of
isomorphisms, formed using isomorphisms of the form $c_{u,v}$, is commutative:
\[
\begin{array}{ccc}
  (h(gf))^*=((hg)f)^* & \longrightarrow & f^*(hg)^* \\
  \big\downarrow && \big\downarrow \\
  (gf)^*h^* & \longrightarrow & (f^*g^*)h^*=f^*(g^*h^*).
\end{array}
\]
\end{enumerate}

\result{EXAMPLE}{1} Let $\mathcal C$ be a category in which fiber products
exist. A fibered category $\mathcal F$ over $\mathcal C$ is then defined by
letting $\mathcal F_X$ be the category of objects of $\mathcal C$ over $X$;
the functor $f^*:\mathcal F_Y\longrightarrow\mathcal F_X$ corresponding to a
morphism $f:X\longrightarrow Y$ is defined by the fiber product
$Z\longmapsto Z\times_YX$.

\result{EXAMPLE}{2} Let $\mathcal C$ be the category of preschemes, and for
$X\in\mathcal C$ let $\mathcal F_X$ be the category of quasi-coherent sheaves
of modules on $X$. If $f:X\longrightarrow Y$ is a morphism of preschemes, then
$f^*:\mathcal F_Y\longrightarrow\mathcal F_X$ is the inverse-image functor

\src{301}
on sheaves of modules. This gives a fibered category over $\mathcal C$.

\result{b. DEFINITION}{1.2} A diagram
\[
  E\xrightarrow{u}E'
  \overset{v_1}{\underset{v_2}{\rightrightarrows}}E''
\]
of maps of sets is called \emph{exact} if $u$ is a bijection from $E$ onto the
part of $E'$ consisting of the $x'\in E'$ such that $v_1(x')=v_2(x')$.

\result{DEFINITION}{1.3} Let $\mathcal F$ be a fibered category over
$\mathcal C$, and consider a diagram
\[
  \tag{+}
  S\xleftarrow{\alpha}S'
  \overset{\beta_1}{\underset{\beta_2}{\rightrightarrows}}S''
\]
of morphisms in $\mathcal C$ such that
$\alpha\beta_1=\alpha\beta_2$. This diagram is called
$\mathcal F$-\emph{exact} if, for every pair $\xi,\eta$ of objects of
$\mathcal F_S$, the following diagram of maps of sets is exact:
\[
  \tag{+'}
  \operatorname{Hom}(\xi,\eta)
  \xrightarrow{\alpha^*}
  \operatorname{Hom}(\alpha^*\xi,\alpha^*\eta)
  \overset{\beta_1^*}{\underset{\beta_2^*}{\rightrightarrows}}
  \operatorname{Hom}(\gamma^*\xi,\gamma^*\eta),
\]
where $\gamma=\alpha\beta_1=\alpha\beta_2$. In the latter diagram we have
identified $\beta_i^*\alpha^*$ with
$(\alpha\beta_i)^*=\gamma^*$ by means of $c_{\beta_i,\alpha}$, for simplicity.

\result{DEFINITION}{1.4} Let $\mathcal F$ be a fibered category over
$\mathcal C$, and consider two morphisms
$\beta_1,\beta_2:S''\longrightarrow S'$ in $\mathcal C$. Let
$\xi'\in\mathcal F_{S'}$. A \emph{gluing datum} on $\xi'$ (relative to the pair
$(\beta_1,\beta_2)$) is an isomorphism from $\beta_1^*(\xi')$ to
$\beta_2^*(\xi')$. If $\xi',\eta'\in\mathcal F_{S'}$ are each equipped with a
gluing datum, a morphism $u:\xi'\longrightarrow\eta'$ in
$\mathcal F_{S'}$ is said to be \emph{compatible with the gluing data} if the
following diagram is commutative:
\[
\begin{array}{ccc}
  \beta_1^*(\xi') & \longrightarrow & \beta_2^*(\xi') \\
  \big\downarrow && \big\downarrow \\
  \beta_1^*(\eta') & \longrightarrow & \beta_2^*(\eta').
\end{array}
\]
In this way the objects of $\mathcal F_{S'}$ equipped with a gluing datum
(relative to $\beta_1,\beta_2$) form a category. If
$\alpha:S'\longrightarrow S$ is a morphism such that
$\alpha\beta_1=\alpha\beta_2$, then for every $\xi\in\mathcal F_S$ the object
$\xi'=\alpha^*(\xi)$

\src{302}
of $\mathcal F_{S'}$ carries a canonical gluing datum, since
\[
  \beta_i^*\alpha^*(\xi)\simeq(\alpha\beta_i)^*(\xi)=\gamma^*(\xi),
\]
where again $\gamma=\alpha\beta_1=\alpha\beta_2$; moreover, if
$u:\xi\longrightarrow\eta$ is a morphism in $\mathcal F_S$, then
\[
  \alpha^*(u):\alpha^*(\xi)\longrightarrow\alpha^*(\eta)
\]
is a morphism in $\mathcal F_{S'}$ compatible with the canonical gluing data.
Thus one obtains a canonical functor from $\mathcal F_S$ to the category of
objects of $\mathcal F_{S'}$ equipped with a gluing datum relative to the pair
$(\beta_1,\beta_2)$. Definition 1.3 may therefore also be expressed by saying
that the diagram (+) there is $\mathcal F$-exact if the preceding functor is
“fully faithful,” i.e. defines an equivalence of $\mathcal F_S$ with a full
subcategory of the category of objects of $\mathcal F_{S'}$ equipped with a
gluing datum relative to $(\beta_1,\beta_2)$.

\result{DEFINITION}{1.5} A gluing datum on
$\xi'\in\mathcal F_{S'}$ is called \emph{effective} (relative to $\alpha$) if
$\xi'$, equipped with this datum, is isomorphic to some $\alpha^*(\xi)$ with
$\xi\in\mathcal F_S$.

When the diagram (+) is $\mathcal F$-exact, the object $\xi$ in the preceding
definition is therefore determined up to unique isomorphism, and the category
$\mathcal F_S$ is equivalent to the category of objects of
$\mathcal F_{S'}$ equipped with an effective gluing datum.

\medskip\noindent
c. The most important case is that in which
\[
  S''=S'\times_S S',
\]
the $\beta_i$ being the two projections $p_1$ and $p_2$ from
$S'\times_S S'$ onto its two factors (we now assume that fiber products exist
in $\mathcal C$). There are then two immediately necessary conditions for a
gluing datum $\varphi:p_1^*(\xi')\longrightarrow p_2^*(\xi')$ on an object
$\xi'\in\mathcal F_{S'}$ to be effective. One has
\[
  \tag{i}
  \Delta^*(\varphi)=\operatorname{id}_{\xi'},
\]
where $\Delta:S'\longrightarrow S'\times_S S'$ denotes the diagonal morphism,
and where $\Delta^*p_i^*(\xi')$ has been identified with
$(p_i\Delta)^*(\xi')=\xi'$. One also has
\[
  \tag{ii}
  p_{31}^*(\varphi)=p_{32}^*(\varphi)\,p_{21}^*(\varphi),
\]

\src{303}
where $p_{ij}$ denotes the canonical projection from
$S'\times_S S'\times_S S'$ onto the partial product of its $i$th and $j$th
factors.

\result{DEFINITION}{1.6} A \emph{descent datum} on
$\xi'\in\mathcal F_{S'}$, relative to the morphism
$\alpha:S'\longrightarrow S$, is a gluing datum on $\xi'$ relative to the
pair $(p_1,p_2)$ of canonical projections
$S'\times_S S'\longrightarrow S'$ satisfying conditions (i) and (ii) above.

\result{DEFINITION}{1.7} A morphism $\alpha:S'\longrightarrow S$ is called an
\emph{$\mathcal F$-descent morphism} if the diagram
\[
  S\xleftarrow{\alpha}S'
  \overset{p_1}{\underset{p_2}{\rightrightarrows}}S'\times_S S'
\]
is $\mathcal F$-exact (Definition 1.3); one says that $\alpha$ is a
\emph{strict $\mathcal F$-descent morphism} if, in addition, every descent
datum (Definition 1.6) on an object of $\mathcal F_{S'}$ is effective.

The latter condition may also be stated more pictorially as follows: “it
amounts to the same thing” to specify an object of $\mathcal F_S$ or an object
of $\mathcal F_{S'}$ equipped with a descent datum.

We note that if a morphism $\alpha:S'\longrightarrow S$ admits a “section”
$s:S\longrightarrow S'$ (i.e. a morphism $s$ such that
$\alpha s=\operatorname{id}_S$), then it is a strict
$\mathcal F$-descent morphism.\footnote{The 1962 erratum specifies that it is
unnecessary to assume beforehand that $\alpha$ is an $\mathcal F$-descent
morphism.} Indeed, if $\xi'\in\mathcal F_{S'}$ is equipped with a descent datum
relative to $\alpha$, it “comes from” $\xi=s^*(\xi')$.

\medskip\noindent
d. The preceding notions may be presented more intuitively by introducing, for
a parameter object $T$ of $\mathcal C$ over $S$, the set
\[
  \operatorname{Hom}_S(T,S')=S'(T),
\]
whose elements will be denoted by $t,t'$, and so forth. Given an object
$\xi'\in\mathcal F_{S'}$, each $t\in S'(T)$ then determines an object
$t^*(\xi')$ of $\mathcal F_T$, also denoted by $\xi'_t$. A gluing datum on
$\xi'$ (relative to $p_1,p_2$) is then specified, for every $T$ over $S$ and
every pair of points $t,t'\in S'(T)$, by an isomorphism
\[
  \varphi_{t',t}:\xi'_t\longrightarrow\xi'_{t'}
\]
(satisfying the evident functoriality conditions as $T$ varies). Conditions (i)
and (ii) of 1.c then become
\[
  \tag{i bis}
  \varphi_{t,t}=\operatorname{id}_{\xi'_t}
  \qquad\text{for every $T$ and every $t\in S'(T)$},
\]

\src{304}
\[
  \tag{ii bis}
  \varphi_{t,t''}=\varphi_{t,t'}\varphi_{t',t''}
  \qquad\text{for every $T$ and all $t,t',t''\in S'(T)$}.
\]
It is clear moreover that (ii bis), on setting $t=t'=t''$, implies
$\varphi_{t,t}=\varphi_{t,t}^2$; since $\varphi_{t,t}$ is an isomorphism by
hypothesis, this gives (i bis), which is therefore actually a consequence of
(ii bis) (and hence (i) is a consequence of (ii)). But if the
$\varphi_{t,t'}$ are no longer assumed a priori to be isomorphisms (i.e. if
$\varphi:p_1^*(\xi')\longrightarrow p_2^*(\xi')$ is not assumed to be an
isomorphism), then (ii bis) no longer necessarily implies (i bis); the
conjunction of (ii bis) and (i bis) nevertheless implies that the
$\varphi_{t,t'}$ are isomorphisms (for then
$\varphi_{t,t'}\varphi_{t',t}=\varphi_{t,t}=\operatorname{id}_{\xi'_t}$).

\subsection*{2. Exact diagrams and strict epimorphisms, descent morphisms.
Examples}
\addcontentsline{toc}{subsection}{2. Exact diagrams and strict epimorphisms}

\result{a. DEFINITION}{2.1} Let $\mathcal C$ be a category. A diagram
\[
  T\xrightarrow{\alpha}T'
  \overset{\beta_1}{\underset{\beta_2}{\rightrightarrows}}T''
\]
of morphisms in $\mathcal C$ is called \emph{exact} if, for every
$Z\in\mathcal C$, the corresponding diagram of maps of sets
\[
  \operatorname{Hom}(Z,T)\longrightarrow
  \operatorname{Hom}(Z,T')
  \rightrightarrows\operatorname{Hom}(Z,T'')
\]
is exact (Definition 1.2). One then says that $(T,\alpha)$ (or, by abuse of
language, that $T$) is a \emph{kernel} of the pair of morphisms
$(\beta_1,\beta_2)$.

This kernel is plainly determined up to unique isomorphism. If $\mathcal C$ is
the category of sets, the preceding definition is compatible with Definition
1.2. Dually, one defines exactness of a diagram
\[
  S\xleftarrow{\alpha}S'
  \overset{\beta_1}{\underset{\beta_2}{\rightrightarrows}}S''
\]
of morphisms in $\mathcal C$; one then says that $(S,\alpha)$ is a
\emph{cokernel} of the pair of morphisms $(\beta_1,\beta_2)$.

\result{DEFINITION}{2.2} A morphism $\alpha:S'\longrightarrow S$ is called a
\emph{strict epimorphism} if it is an epimorphism and if, for every morphism
$u:S'\longrightarrow Z$, the following necessary condition is also sufficient
for $u$ to factor as $S'\longrightarrow S\longrightarrow Z$: for every
$S''\in\mathcal C$ and every pair of morphisms
$\beta_1,\beta_2:S''\longrightarrow S'$ such that
$\alpha\beta_1=\alpha\beta_2$, one also has $u\beta_1=u\beta_2$.

If the fiber product $S'\times_S S'$ exists, this is equivalent to saying that
the diagram
\[
  S\xleftarrow{\alpha}S'
  \overset{p_1}{\underset{p_2}{\rightrightarrows}}S'\times_S S'
\]

\src{305}
is exact, i.e. $S$ is a cokernel of the pair $(p_1,p_2)$. In any case, a
cokernel morphism is a strict epimorphism. Note also that a strict epimorphism
which is a monomorphism is an isomorphism. We leave it to the reader to develop
the dual notion of a \emph{strict monomorphism}.

To make precise the relations between the notion of an $\mathcal F$-descent
morphism and the notion of a strict epimorphism, we introduce the following
further definitions.

\result{DEFINITION}{2.3} A morphism $\alpha:S'\longrightarrow S$ is called a
\emph{universal epimorphism} (resp. a \emph{universal strict epimorphism}) if,
for every $T$ over $S$, the fiber product $T'=S'\times_ST$ exists and the
projection $T'\longrightarrow T$ is an epimorphism (resp. a strict
epimorphism).

In very good categories (such as the category of sets, the category of sheaves
of sets on a topological space, abelian categories, and so forth), the four
notions of epimorphism just introduced coincide. By contrast, they are all
distinct in a category such as the category of preschemes, or the category of
preschemes over a given nonempty prescheme $S$, even if one restricts to
$S$-schemes finite over $S$.

\result{DEFINITION}{2.4} A morphism $\alpha:S'\longrightarrow S$ is called a
\emph{descent morphism} (resp. a \emph{strict descent morphism}) if it is an
$\mathcal F$-descent morphism (resp. a strict $\mathcal F$-descent morphism)
(cf. Definition 1.7), where $\mathcal F$ denotes the fibered category over
$\mathcal C$ of objects of $\mathcal C$ over objects of $\mathcal C$ (no. 1,
Example 1).

\result{PROPOSITION}{2.1} If finite products and (finite) fiber products exist
in $\mathcal C$, then in $\mathcal C$ descent morphisms and universal strict
epimorphisms are the same.

\result{b. EXAMPLES}{} Let $\mathcal C$ be the category of preschemes. Let
$S\in\mathcal C$, and let $S'$ and $S''$ be two preschemes finite over $S$,
i.e. corresponding to sheaves of algebras $\mathcal A'$, $\mathcal A''$ on $S$
which, as sheaves of modules, are quasi-coherent and of finite type (i.e.
coherent if $S$ is locally Noetherian). Let $\alpha:S'\longrightarrow S$ be the
structure morphism of $S'$, and let $\beta_1,\beta_2$ be two $S$-morphisms from
$S''$ to $S'$, defined by homomorphisms of algebras
$\mathcal A'\longrightarrow\mathcal A''$, again denoted by $\beta_1,\beta_2$.
Using the fact that a finite morphism is closed (the first Cohen--Seidenberg
theorem), one proves easily that the diagram
\[
  \tag{+}
  S\xleftarrow{\alpha}S'
  \overset{\beta_1}{\underset{\beta_2}{\rightrightarrows}}S''
\]
in $\mathcal C$ is exact if and only if the diagram of sheaves

\src{306}
\[
  \mathcal O_S=\mathcal A\longrightarrow\mathcal A'
  \overset{\beta_1}{\underset{\beta_2}{\rightrightarrows}}\mathcal A''
\]
on $S$ is exact. In particular, if $\alpha:S'\longrightarrow S$ is a finite
morphism corresponding to a sheaf $\mathcal A'$ of algebras on $S$, then
$\alpha$ is a strict epimorphism if and only if the diagram of sheaves
\[
  \mathcal O_S=\mathcal A\longrightarrow\mathcal A'
  \overset{p_1}{\underset{p_2}{\rightrightarrows}}
  \mathcal A'\otimes_{\mathcal A}\mathcal A'
\]
is exact (it is an epimorphism if and only if
$\mathcal A\longrightarrow\mathcal A'$ is injective). If $S$ is affine with
ring $A$, and hence $S'$ is affine with ring $A'$ finite over $A$, then
$S'\longrightarrow S$ is a strict epimorphism if and only if
$A\longrightarrow A'$ is an isomorphism from $A$ onto the subring of $A'$
formed by those $x'\in A'$ such that
\[
  1_{A'}\otimes_Ax'-x'\otimes_A1_{A'}=0
\]
(it is an epimorphism if and only if $A\longrightarrow A'$ is injective). As
we have already pointed out, even if $S$ is the scheme of a local Artinian
ring, a finite morphism $S'\longrightarrow S$ which is an epimorphism need not
be a strict epimorphism. It can nevertheless be proved that, if $S$ is a
Noetherian prescheme, every finite morphism $S'\longrightarrow S$ which is an
epimorphism is the composite of a finite sequence of strict epimorphisms (also
finite). This also shows that the composite of two strict epimorphisms need not
be a strict epimorphism.

\result{c}{} If $(+)$ is an exact diagram of finite morphisms, then, for every
flat morphism $T\longrightarrow S$ of preschemes, the diagram obtained from
$(+)$ by the base change $T\longrightarrow S$ is again exact. It follows that,
if $X,Y$ are two $S$-preschemes, with $X$ flat over $S$, then the following
diagram of maps of sets (where $X',Y'$ are the inverse images of $X,Y$ over
$S'$, and $X'',Y''$ their inverse images over $S''$)
\[
  \operatorname{Hom}_S(X,Y)\longrightarrow
  \operatorname{Hom}_{S'}(X',Y')
  \rightrightarrows
  \operatorname{Hom}_{S''}(X'',Y'')
\]
is exact. In particular, if $\mathcal F$ denotes the fibered category over the
category $\mathcal C$ of preschemes such that, for $X\in\mathcal C$,
$\mathcal F_X$ is the category of flat $X$-preschemes, then the diagram $(+)$
is $\mathcal F$-exact. (This result becomes false if no flatness hypothesis is
made; in particular, a finite strict epimorphism need not be a descent
morphism.) Likewise, $(+)$ is $\mathcal F$-exact if $\mathcal F$ denotes the
fibered category for which $\mathcal F_X$ is the category of quasi-coherent flat
sheaves on the prescheme $X$ (here again the flatness hypothesis

\src{307}
is essential). In either case, the question of the effectivity of a gluing
datum (and more particularly of a descent datum when $S''=S'\times_SS'$) on an
object flat over $S'$ is delicate (and its answer in various special cases is
one of the main subjects of the present talks). The speaker does not know
whether, for every finite strict epimorphism $S'\longrightarrow S$, every
descent datum on a flat quasi-coherent sheaf on $S'$ is effective (even if $S$
is assumed to be the spectrum of a local Artinian ring and one restricts to
locally free sheaves of rank 1). More generally, let $A$ be a ring and $A'$ an
$A$-algebra (all rings are commutative) such that the diagram of maps
\[
  A\longrightarrow A'\rightrightarrows A'\otimes_AA'
\]
is exact; equivalently, the corresponding morphism $S'\longrightarrow S$ on
the spectra of $A,A'$ is an $\mathcal F$-descent morphism, where $\mathcal F$
is the fibered category of flat quasi-coherent sheaves. Let $M'$ be a flat
$A'$-module equipped with a descent datum to $A$, i.e. an isomorphism
\[
  \varphi:M'\otimes_AA'\xrightarrow{\sim}A'\otimes_AM'
\]
of $A'\otimes_AA'$-modules satisfying conditions (i) and (ii) of 1.c (whose
formulation in terms of modules we leave to the reader). Is this datum
effective (relative to the fibered category of flat quasi-coherent sheaves)?
Let $M$ be the subset of $M'$ formed by the $x'\in M'$ such that
\[
  \varphi(x'\otimes_A1_{A'})=1_{A'}\otimes_Ax'.
\]
It is an $A$-submodule of $M'$. The canonical injection
$M\longrightarrow M'$ defines a homomorphism of $A'$-modules
$M\otimes_AA'\longrightarrow M'$. The effectivity of $\varphi$ then means the
following: $M$ is a flat $A$-module, and the preceding homomorphism is an
isomorphism.

\result{REMARK}{} In the preceding considerations we made no flatness
hypothesis on the morphisms in the diagram $(+)$; to obtain a descent
technique, this compelled us to make flatness hypotheses on the objects over
$S,S'$ under consideration. In Section 2 we shall make a flatness hypothesis
on $\alpha:S'\longrightarrow S$, which will allow us to have a descent
technique for objects over $S,S'$ that are no longer subject to any flatness
condition. In every case, a flatness hypothesis enters. This is one of the
principal reasons for the importance of the notion of flatness in algebraic
geometry (a role which could not become apparent so long as one restricted to
base fields, over which anything whatever is, indeed, flat!).

\src{308}
\subsection*{3. Application to étale morphisms}
\addcontentsline{toc}{subsection}{3. Application to étale morphisms}

Let $A$ be a local ring and $B$ a local algebra over $A$ whose maximal ideal
induces that of $A$. We shall say that $B$ is \emph{étale} over $A$ (instead of
“unramified,” used elsewhere) if it satisfies the following conditions:
\begin{enumerate}
\item[(i)] $B$ is flat over $A$;
\item[(ii)] $B/\mathfrak mB$ is a finite separable extension of
$A/\mathfrak m=k$ (where $\mathfrak m$ denotes the maximal ideal of $A$).
\end{enumerate}
When $A$ and $B$ are Noetherian and $k$ is algebraically closed, this means that
the homomorphism $\widehat A\longrightarrow\widehat B$ on completions extending
$A\longrightarrow B$ is an isomorphism. A morphism of finite type
$f:T\longrightarrow S$ is said to be \emph{étale at $x\in T$}, or $T$ is said
to be \emph{étale over $S$ at $x$}, if $\mathcal O_x$ is étale over
$\mathcal O_{f(x)}$; and $f$ is said to be \emph{étale}, or $f$ is called an
\emph{étale morphism}, or $T$ is said to be \emph{étale over $S$}, if $f$ is
étale at every $x\in T$. Observe moreover that, if $S$ is locally Noetherian,
the set of points of $T$ where $f$ is étale is open, and the use of Zariski's
“main theorem” makes it possible to specify the structure of $T/S$ in a
neighborhood of such a point (by an equation of a well-known type).

When $S$ is a scheme of finite type over the field of complex numbers, there
corresponds to it an analytic space $\overline S$ in the sense of Serre [5],
except that $\overline S$ may have nilpotent elements in its structure sheaf,
which changes nothing essential in [5]. It is then easy to see that $f$ is
étale if and only if $\overline f:\overline T\longrightarrow\overline S$ is
étale, i.e. if every point of $\overline T$ has a neighborhood on which
$\overline f$ induces an isomorphism onto an open subset of $\overline S$. In
particular, to every \emph{étale covering} $T$ of $S$ (i.e. a finite étale
morphism $f:T\longrightarrow S$) there corresponds an étale covering
$\overline T$ of $\overline S$, which is connected if and only if $T$ is [5].
It is also easy to see that, if $T,T'$ are two schemes étale over $S$, then the
natural map
\[
  \operatorname{Hom}_S(T,T')\longrightarrow
  \operatorname{Hom}_{\overline S}(\overline T,\overline{T'})
\]
is bijective, i.e. the functor $T\longmapsto\overline T$ from the category of
schemes étale over $S$ to the category of analytic spaces étale over
$\overline S$ is “fully faithful,” and therefore defines an equivalence of the
first category with a subcategory of the second. A theorem of
Grauert--Remmert [2] implies that, if $S$ is normal, one thereby obtains an
equivalence between the category of \emph{étale coverings} of $S$ and the
category of (finite) étale coverings of $\overline S$, i.e. every étale
covering of $\overline S$ is $\overline S$-isomorphic to some $\overline T$,
where $T$ is an étale covering of $S$. Let us show that the
\emph{Grauert--Remmert theorem remains valid without a normality hypothesis on
$S$}. First let $S'\longrightarrow S$ be a finite strict epimorphism, assume
the

\src{309}
theorem proved for $S'$, and let us show that it is true for $S$. Indeed, let
$\mathfrak T$ be an étale covering of $\overline S$, and consider its inverse
image $\mathfrak T'$ over $\overline{S'}$, which corresponds to a coherent
analytic sheaf $\mathcal A'$ of algebras on $\overline{S'}$, the inverse image
of the sheaf of algebras $\mathcal A$ on $\overline S$ defining
$\mathfrak T$. By hypothesis, over $\overline{S'}$, $\mathfrak T'$ comes from
an étale covering $T'$ of $S'$, i.e. $\mathcal A'$ comes from a coherent sheaf
of algebras $A'$ on $S'$. On the other hand, $\mathcal A'$ is equipped with a
canonical descent datum relative to
$\overline{S'}\longrightarrow\overline S$, i.e. an isomorphism between its two
inverse images on
\[
  \overline{S'}\times_{\overline S}\overline{S'}
  =\overline{(S'\times_SS')},
\]
(satisfying conditions (i), (ii)); and, by what was said above, this
isomorphism comes from an isomorphism on the corresponding algebraic sheaves,
i.e. from a descent datum on $A'$ relative to $S'\longrightarrow S$. One checks
easily that the latter is effective (for the datum on $\mathcal A'$ is
effective, and the effectivity of a descent datum, as made explicit in the
preceding subsection, can be recognized locally on the \emph{completions} of
the modules involved). Hence there is a coherent sheaf of algebras $A$ on $S$
defining a covering $T$ of $S$, which is the desired covering. The preceding
result plainly remains valid if $S'\longrightarrow S$ is merely a composite of
a finite number of finite strict epimorphisms, i.e. an arbitrary finite
epimorphism (by the factorization result noted in Section 2). It follows that
the Grauert--Remmert theorem remains valid if $S$ is a \emph{reduced} scheme,
i.e. if $\mathcal O_S$ has no nilpotent elements, as one sees by introducing
its normalization $S'$. The general case follows easily.

A wholly analogous proof, again using the factorization result for finite
strict epimorphisms and the “formal” nature of the effectivity of descent data,
proves the following result: let $S$ be a locally Noetherian prescheme and let
$S'\longrightarrow S$ be a morphism that is \emph{finite, surjective, and
radicial} (or, equivalently, a morphism of finite type such that, for every $T$
over $S$, the morphism $T'=S'\times_ST\longrightarrow T$ is a homeomorphism,
which is also expressed by saying that $S'\longrightarrow S$ is a
\emph{universal homeomorphism}). For every $T$ étale over $S$, consider its
inverse image $T'=T\times_SS'$, which is étale over $S'$. Then the functor
$T\longmapsto T'$ is an \emph{equivalence of categories from the category of
preschemes $T$ étale over $S$ to the category of preschemes $T'$ étale over
$S'$}. (One uses the bijectivity of
\[
  \operatorname{Hom}_S(T_1,T_2)\longrightarrow
  \operatorname{Hom}_{S'}(T'_1,T'_2)
\]
for two preschemes $T_1,T_2$ étale over $S$, a fact whose direct verification
is easy, and the fact that the stated theorem is true if
$S'=(S,\mathcal O_S/\mathcal J)$,

\src{310}
where $\mathcal J$ is a coherent nilpotent sheaf of ideals of $\mathcal O_S$
([4], Lemma 6)). Note moreover that here we do not assume the $T$ under
consideration to be finite over $S$. This result implies in particular that the
morphism $S'\longrightarrow S$ induces an isomorphism from the fundamental
group $[\,]$ of $S'$ onto that of $S$ (“\emph{topological invariance of the
fundamental group of a prescheme}”).

\subsection*{4. Relations with 1-cohomology}
\addcontentsline{toc}{subsection}{4. Relations with 1-cohomology}

\medskip\noindent
a. Let $\mathcal C$ be a category in which the product of any two objects
always exists, and let $T\in\mathcal C$. For every finite set
$I\ne\varnothing$, one can consider $T^I$; as $I$ varies, this yields a
covariant functor from the category of nonempty finite sets to $\mathcal C$,
i.e. what may be called a \emph{simplicial object} of $\mathcal C$, denoted by
$K_T$. The latter depends covariantly on $T$; moreover, if $u,v$ are two
morphisms $T\longrightarrow T'$, then the corresponding morphisms
$K_T\longrightarrow K_{T'}$ are homotopic. Say that $T$ \emph{dominates} $T'$
if $\operatorname{Hom}(T,T')\ne\varnothing$; this is an increasing filtered
preorder relation in $\mathcal C$. It follows from the foregoing that, if $T$
dominates $T'$, there is a canonical homotopy class of homomorphisms of
simplicial objects $K_T\longrightarrow K_{T'}$; in particular, if $K_T$ and
$K_{T'}$ are such that each dominates the other, then $K_T$ and $K_{T'}$ are
homotopy equivalent. Now let $F$ be a functor (contravariant, to fix ideas) from
$\mathcal C$ to an abelian category $\mathcal C'$. Then
\[
  C^*(T,F)=F(K_T)
\]
is a cosimplicial object of $\mathcal C'$, and hence defines in the usual way a
(cochain) complex in $\mathcal C'$, whose cohomology may be taken:
\[
  H^*(T,F)=H^*(C^*(T,F))=H^*(F(K_T)).
\]
(One may put a $\mathcal C$ as a subscript on $H^*$ if there is a possibility
of confusion.) This is a cohomological functor in $F$, whose variance in the
variable $T$ follows from what was said about the $K_T$; more precisely, for
$F$ fixed and $T$ varying in $\mathcal C$ (preordered by the domination
relation), the $H^*(T,F)$ form an inductive system of graded objects of
$\mathcal C'$. In particular, if $T$ and $T'$ are such that each dominates the
other, then $H^*(T,F)$ and $H^*(T',F)$ are canonically isomorphic.

Suppose that fiber products exist in $\mathcal C$. Then, for a fixed
$S\in\mathcal C$, one may apply the foregoing to the category $\mathcal C_S$ of
objects of $\mathcal C$ over $S$; we shall write $C^*(T/S,F)$ and
$H^*(T/S,F)$ instead of $C^*(T,F)$ and $H^*(T,F)$ if we wish to specify that
we are working in the category $\mathcal C_S$; thus,

\src{311}
$C^*(T/S,F)$ is a cochain complex in $\mathcal C'$ whose term in degree $n$ is
\[
  F(T\times_ST\times_S\cdots\times_ST)
\]
(where the parenthesis contains $n+1$ factors).

As usual, note that one may define $H^0(T/S,F)$ without assuming that
$\mathcal C'$ is abelian: it is the kernel (Definition 2.1), if it exists, of
the pair of morphisms $F(p_i)$ ($i=1,2$)
\[
  F(T)\rightrightarrows F(T\times_ST)
\]
corresponding to the two projections
$p_1,p_2:T\times_ST\rightrightarrows T$. In particular, there is a natural
morphism (called the \emph{augmentation})
\[
  F(S)\longrightarrow H^0(T/S,F),
\]
which is an isomorphism in favorable cases (in particular if
$T\longrightarrow S$ is a strict epimorphism and $F$ is “left exact”).
Likewise, when $F$ takes values in the category of groups in a category
$\mathcal C''$, one may also define $H^1(T/S,F)$. When $\mathcal C''$ is the
category of sets (i.e. $F$ takes values in the category of ordinary, not
necessarily commutative, groups), $H^1(T/S,F)$ is the quotient of the subgroup
$Z^1(T/S,F)$ of $C^1(T/S,F)=F(T\times_ST)$ formed by the $g$ such that
\[
  F(p_{31})(g)=F(p_{32})(g)\,F(p_{21})(g)
\]
by the operator group $F(T)$, acting on $C^1(T/S,F)$ and in particular on the
subset $Z^1(T/S,F)$ by
\[
  \rho(g')\cdot g=F(p_2)(g')\,g\,F(p_1)(g')^{-1}.
\]

\medskip\noindent
b. For example, let $\mathcal F$ be a fibered category over $\mathcal C$. Let
$\xi,\eta\in\mathcal F_S$, and for every $S'$ over $S$, set
\[
  F_{\xi,\eta}(S')=
  \operatorname{Hom}(\xi\times_SS',\eta\times_SS').
\]
Thus $F_{\xi,\eta}$ is a contravariant functor from $\mathcal C_S$ to the
category of sets. With this notation, \emph{to say that the augmentation
morphism}
\[
  F_{\xi,\eta}(S)\longrightarrow H^0(S'/S,F_{\xi,\eta})
\]
\emph{is an isomorphism for every pair of elements
$\xi,\eta\in\mathcal F_S$ means that
$\alpha:S'\longrightarrow S$ is an $\mathcal F$-descent morphism}
(Definition 1.7).

\medskip\noindent
c. Similarly, for $\xi\in\mathcal F_S$ and every object $S'$ of $\mathcal C$
over $S$, set
\[
  G_\xi(S')=\operatorname{Aut}(\xi\times_SS').
\]
This defines a contravariant functor $G_\xi$ from $\mathcal C_S$ to the
category of

\src{312}
groups. One then sees that $Z^1(S'/S,G_\xi)$ \emph{is canonically identified
with the set of descent data on $\xi'=\xi\times_SS'$ relative to
$S'\longrightarrow S$} (Definition 1.6), and $H^1(S'/S,G_\xi)$ \emph{is
identified with the set of classes, up to isomorphism, of objects of
$\mathcal F_{S'}$ equipped with a descent datum relative to
$\alpha:S'\longrightarrow S$ which, as objects of $\mathcal F_{S'}$, are
isomorphic to $\xi'=\xi\times_SS'$}. Thus, if
$\alpha:S'\longrightarrow S$ is an $\mathcal F$-descent morphism (cf. (b)),
then $H^1(S'/S,G_\xi)$ \emph{contains as a subset the set of classes, up to
isomorphism, of objects $\eta$ of $\mathcal F_S$ such that
$\eta\times_SS'$ is isomorphic in $\mathcal F_{S'}$ to $\xi\times_SS'$; and
this inclusion is an identity if and only if every descent datum on
$\xi'=\xi\times_SS'$ relative to $\alpha:S'\longrightarrow S$ is effective}.
(This will hold in particular if $\alpha:S'\longrightarrow S$ is a strict
$\mathcal F$-descent morphism.)

\result{REMARK}{} Cochain complexes of the form $C^*(T/S,F)$ contain as special
cases most of the familiar standard complexes (Čech cohomology, group
cohomology, and so forth), and play an important role in algebraic geometry,
notably in the “Weil cohomology” of preschemes.

\medskip\noindent
d. \result{EXAMPLE}{1} Let $S'$ be an object over $S\in\mathcal C$, and let
$\Gamma$ be a group of automorphisms of $S'$ such that $S'$ is “formally
principal over $S$, with group $\Gamma$,” i.e. such that the natural morphism
\[
  \Gamma\times S'\longrightarrow S'\times_SS',
\]
where $\Gamma\times S'$ denotes the coproduct of $\Gamma$ copies of $S'$, is an
isomorphism. (We assume that the coproducts occurring here exist in
$\mathcal C$.) Let $F$ be a contravariant functor from $\mathcal C$ to the
category of abelian groups. Then $C^*(S'/S,F)$ is canonically isomorphic to the
simplicial group of standard homogeneous cochains $C^*(\Gamma,F(S'))$; hence
$H^*(S'/S,F)$ is canonically isomorphic to $H^*(\Gamma,F(S'))$.

\medskip\noindent
e. \result{EXAMPLE}{2} Let $\mathcal C$ be the category of preschemes. Denote by
$G_a$ (“additive group”) the contravariant functor from $\mathcal C$ to the
category of abelian groups defined by
\[
  G_a(X)=H^0(X,\mathcal O_X).
\]
Similarly, define the functor $G_m$ (“multiplicative group”) by
\[
  G_m(X)=H^0(X,\mathcal O_X)^*
\]
(the group of invertible elements of the ring $H^0(X,\mathcal O_X)$), and more
generally the functor $\operatorname{Gl}(n)$ (“linear group of order $n$”) by

\src{313}
\[
  \operatorname{Gl}(n)(X)=
  \operatorname{Gl}\bigl(n,H^0(X,\mathcal O_X)\bigr).
\]
This is a functor from $\mathcal C$ to the category of groups (not necessarily
commutative if $n>1$; for $n=1$ it is $G_m$). Moreover,
$\operatorname{Gl}(n)$ may be interpreted as an automorphism functor (cf. (c))
by considering the fibered category $\mathcal F$ over $\mathcal C$ such that,
for $X\in\mathcal C$, $\mathcal F_X$ is the category of locally free sheaves on
$X$: indeed,
\[
  \operatorname{Gl}(n)(X)=
  \operatorname{Aut}_{\mathcal F_X}(\mathcal O_X^n).
\]
By (b), if $\alpha:S'\longrightarrow S$ is an
$\mathcal F$-descent morphism (cf. Section 2.c),
$H^1(S'/S,\operatorname{Gl}(n))$ contains the set of classes, up to
isomorphism, of locally free sheaves on $S$ whose inverse image on $S'$ is
isomorphic to $\mathcal O_{S'}^n$; and this inclusion is an equality if and
only if every descent datum on $\mathcal O_{S'}^n$ (relative to
$\alpha:S'\longrightarrow S$) is effective. When $S$ is the spectrum of a
local ring, this therefore means that
$H^1(S'/S,\operatorname{Gl}(n))=(e)$, since every locally free sheaf on $S$ is
then trivial.

Note the equivalence of the following conditions on a morphism
$\alpha:S'\longrightarrow S$:
\begin{enumerate}
\item[(i)] the augmentation homomorphism
\[
  H^0(S,\mathcal O_S)=G_a(S)\longrightarrow H^0(S'/S,G_a)
\]
is an isomorphism;
\item[(ii)] $S'\longrightarrow S$ is an $\mathcal F$-descent morphism
($\mathcal F$ being the fibered category over $\mathcal C$ considered above).
\end{enumerate}
If $S'\longrightarrow S$ is finite, these conditions are also equivalent to:
\begin{enumerate}
\item[(iii)] $S'\longrightarrow S$ is a strict epimorphism
(cf. Section 2.c).
\end{enumerate}

Now suppose that $S=\operatorname{Spec}(A)$ and
$S'=\operatorname{Spec}(A')$. Then
\[
  C^n(S'/S,G_a)=C^n(A'/A,G_a)=\bigotimes_A^{n+1}A',
\]
the coboundary operator
$C^n(A'/A,G_a)\longrightarrow C^{n+1}(A'/A,G_a)$ being the alternating sum of
the face operators
\[
  \partial_i(x_0\otimes x_1\otimes\cdots\otimes x_n)
  =x_0\otimes\cdots\otimes x_{i-1}\otimes1_{A'}\otimes
   x_i\otimes\cdots\otimes x_n.
\]
Likewise,
\[
  C^n(S'/S,G_m)=C^n(A'/A,G_m)
  \simeq\left(\bigotimes_A^{n+1}A'\right)^*,
\]
the simplicial operations in $C^*(A'/A,G_m)$ being induced by those of
$C^*(S'/S,G_a)$. The simplicial operations in
$C^*(A'/A,\operatorname{Gl}(n))$ are made explicit in the same way.
\emph{In all cases known to the speaker,}
\[
  H^i(A'/A,G_a)=0\qquad\text{for $i>0$},
\]
\emph{and if $A$ is local,}
\[
  H^1(A'/A,G_m)=0
  \quad\text{and more generally}\quad
  H^1(A'/A,\operatorname{Gl}(n))=(e)
\]
(when $S'\longrightarrow S$ is an $\mathcal F$-descent morphism, i.e. the
diagram
$A\longrightarrow A'\rightrightarrows A'\otimes_AA'$ is exact; compare
Section 2.c). Note that \emph{Hilbert's “Theorem 90” is nothing other than the
relation}
\src{314}
\[
  H^1(S'/S,G_m)=0
\]
\emph{when $A$ is a field and $A'$ is a finite Galois extension of $A$
(cf. Example 1), and may also be expressed by saying that, in the case under
consideration, $S'\longrightarrow S$ is a strict descent morphism for the
fibered category of locally free sheaves of rank 1.} It is in this latter form
that Hilbert's theorem should be generalized, by varying the hypotheses both
on the morphism $S'\longrightarrow S$ and on the quasi-coherent sheaves under
consideration.

Finally, note the equivalence of the following properties when $A$ is an
Artinian local ring with maximal ideal $\mathfrak m$ and $A'$ is an
$A$-algebra (where, for every integer $k>0$, $A_k$ (respectively $A'_k$)
denotes the ring $A/\mathfrak m^{k+1}$ (respectively
$A'/\mathfrak m^{k+1}A'$)):
\begin{enumerate}
\item[(i)] $H^1(A'_k/A_k,G_a)=0$ for every $k$.
\item[(ii)] $H^1(A'_k/A_k,G_m)=0$ for every $k$.
\item[(iii)] $H^1(A'_k/A_k,\operatorname{Gl}(n))=(e)$ for every $k$ and every
$n$.
\end{enumerate}
If $S'\longrightarrow S$ is a strict epimorphism, the preceding conditions
even imply that it is a strict descent morphism for free modules (of finite
type or not) over $A'$.

\result{REMARK}{} The definition of the groups $H^i(S'/S,G_m)$, when $S$ and
$S'$ are the schemes of fields $A$ and $A'$, is due to Amitsur. The group
$H^2(S'/S,G_m)$ is especially interesting as a “global” variant of the Brauer
group; for this variant one may refer to [1], Chapter VII.

\fsection{B}{Descent by faithfully flat morphisms}

\subsection*{1. Statement of the descent theorems}
\addcontentsline{toc}{subsection}{1. Statement of the descent theorems}

\result{DEFINITION}{1.1} A morphism
$\alpha:S'\longrightarrow S$ of preschemes is called \emph{flat} if, for every
$x'\in S'$, $\mathcal O_{x'}$ is a flat module over the ring
$\mathcal O_{\alpha(x')}$ (i.e.
$\mathcal O_{x'}\otimes_{\mathcal O_{\alpha(x')}}M$ is an exact functor of the
$\mathcal O_{\alpha(x')}$-module $M$). A morphism is called
\emph{faithfully flat} if it is flat and surjective.

For example, if $S=\operatorname{Spec}(A)$ and
$S'=\operatorname{Spec}(A')$, then $S'$ is flat over $S$ if and only if $A'$
is a flat $A$-module, and $S'$ is faithfully flat over $S$ if and only if $A'$
is a faithfully flat $A$-module (i.e. the functor $A'\otimes_A M$ of the
$A$-module $M$ is exact and faithful); in Serre's terminology [5], this also
means that the pair $(A,A')$ is flat. If $S'$ is faithfully flat over $S$,
then the inverse-image functor on quasi-coherent sheaves over $S$ is exact and
faithful. In other words, for a sequence of homomorphisms of quasi-coherent
sheaves

\src{315}
over $S$ to be exact, it is necessary and sufficient that its inverse image
over $S'$ be exact (in particular, for a homomorphism of quasi-coherent sheaves
over $S$ to be a monomorphism, respectively an epimorphism, respectively an
isomorphism, it is necessary and sufficient that its inverse image over $S'$
be so). This property remains true after restriction over any open subset of
$S$, and in this form characterizes faithfully flat morphisms.

\result{DEFINITION}{1.2} A morphism $\alpha:S'\longrightarrow S$ is called
\emph{quasi-compact} if the inverse image of every quasi-compact open subset
$U$ of $S$ is quasi-compact (i.e. a finite union of affine open subsets).

It plainly suffices to check this property for the affine open subsets of
$S$. For example, an affine morphism (i.e. one for which the inverse image of
an affine open subset is affine) is quasi-compact.

The class of flat, respectively faithfully flat, respectively quasi-compact
morphisms is stable under composition and “base extension,” and contains the
isomorphisms.

\result{THEOREM}{1} Let $\alpha:S'\longrightarrow S$ be a faithfully flat and
quasi-compact morphism of preschemes. Then $\alpha$ is a \emph{strict descent
morphism} (A, Definition 1.7) for the fibered category $\mathcal F$ of
quasi-coherent sheaves (A, Section 1, Example 2).

This statement means two things:
\begin{enumerate}
\item[(i)] If $F$ and $G$ are two quasi-coherent sheaves over $S$, and $F'$ and
$G'$ are their inverse images over $S'$, then the natural homomorphism
\[
  \operatorname{Hom}(F,G)\longrightarrow\operatorname{Hom}(F',G')
\]
is a bijection from the first term onto the subgroup of the second formed by
the homomorphisms $F'\longrightarrow G'$ compatible with the canonical descent
data on these sheaves, i.e. those whose inverse images under the two
projections from $S''=S'\times_SS'$ to $S'$ give the same homomorphism
$F''\longrightarrow G''$.
\item[(ii)] Every quasi-coherent sheaf $F'$ over $S'$ equipped with a descent
datum relative to the morphism $\alpha:S'\longrightarrow S$
(A, Definition 1.6) is isomorphic (together with this datum) to the inverse
image of a quasi-coherent sheaf $F$ over $S$.
\end{enumerate}
Taking $F=\mathcal O_S$ in (i), one obtains:

\result{COROLLARY}{1} Let $G$ be a quasi-coherent sheaf over $S$, let $G'$ and
$G''$ be its inverse images over $S'$ and over $S''=S'\times_SS'$, and let
$p_1,p_2$ be the two projections from $S''$ to $S'$. Then the following
diagram of maps of sets

\src{316}
\[
  \Gamma(G)\xrightarrow{\alpha^*}\Gamma(G')
  \underset{p_2^*}{\overset{p_1^*}{\rightrightarrows}}\Gamma(G'')
\]
is exact (A, Definition 1.(a)).

On the other hand, combining (i), (ii), and Definition 1.1 gives:

\result{COROLLARY}{2} Let $G$ be as in Corollary 1. Then there is a one-to-one
correspondence between the quasi-coherent subsheaves of $G$ and the
quasi-coherent subsheaves of $G'$ whose inverse images over $S''$ under the two
projections $p_1,p_2$ give the same subsheaf of $G''$.

There is, of course, an equivalent statement in terms of quotient sheaves. As
seen in A, Section 4(e), Theorem 1 should be regarded as a generalization of
Hilbert's “Theorem 90,” and includes as special cases various formulations in
terms of 1-cohomology. For the proof, one readily reduces to the case
$S=\operatorname{Spec}(A)$, $S'=\operatorname{Spec}(A')$; and for (i) one
readily reduces to proving Corollary 1, i.e. the exactness of the diagram
\[
  M=A\otimes_A M\longrightarrow A'\otimes_A M
  \rightrightarrows A'\otimes_AA'\otimes_A M
\]
for every $A$-module $M$. This follows from the more general lemma:

\result{LEMMA}{1.1} Let $A'$ be a faithfully flat $A$-algebra. Then, for every
$A$-module $M$, the $M$-augmented complex
$C^*(A'/A,G_a)\otimes_A M$ (cf. A, Section 4(e)) is a \emph{resolution} of
$M$.

It suffices to prove that the augmented complex obtained from the preceding
one by base extension from $A$ to $A'$ satisfies the same conclusions. This
amounts to checking the statement after replacing $A$ by $A'$ and $A'$ by
$A'\otimes_AA'$, and hence reduces us to a case in which there is an
$A$-algebra homomorphism $A'\longrightarrow A$ (or, in geometric terms, to the
case in which $S'$ over $S$ admits a section). In that case the assertion
follows from the general facts of A, Section 4(a). We note in passing the
following corollary, which generalizes a familiar statement in Galois
cohomology (compare A, Section 4(e)):

\result{COROLLARY}{} If $A'$ is faithfully flat over $A$, then
$H^0(A'/A,G_a)=A$ and
\[
  H^i(A'/A,G_a)=0\qquad\text{for $i\geq1$}.
\]

To prove part (ii) of Theorem 1, one proceeds as for (i), reducing to the case
where $S'$ over $S$ admits a section; there it follows from (i)
(cf. A, Section 1(c)).

Theorem 1 and its corollaries may plainly be varied ad libitum by introducing
various additional structures on the quasi-coherent sheaves (or systems of
sheaves) under consideration. For example, the datum over $S$ of a sheaf
\src{317}
of commutative quasi-coherent algebras “is equivalent” to the datum over $S'$
of such a sheaf equipped with a descent datum relative to
$\alpha:S'\longrightarrow S$. In view of the functorial correspondence
between such quasi-coherent sheaves over $S$ and affine preschemes over $S$,
one obtains the second assertion of the following theorem:

\result{THEOREM}{2} Let $\alpha:S'\longrightarrow S$ be as in Theorem 1. Then
$\alpha$ is a \emph{descent morphism} (not strict in general)
(A, Definition 2.4), and it is a \emph{strict descent morphism} for the fibered
category of schemes affine over preschemes (A, Definition 1.7).

The first assertion of the theorem means the following. Let $X,Y$ be two
preschemes over $S$, let $X',Y'$ be their inverse images over $S'$, and let
$X'',Y''$ be their inverse images over $S''=S'\times_SS'$. Then the following
diagram of natural maps
\[
  \operatorname{Hom}_S(X,Y)\xrightarrow{\alpha^*}
  \operatorname{Hom}_{S'}(X',Y')
  \underset{p_2^*}{\overset{p_1^*}{\rightrightarrows}}
  \operatorname{Hom}_{S''}(X'',Y'')
\]
is exact; i.e. $\alpha^*$ is a bijection from
$\operatorname{Hom}_S(X,Y)$ onto the subset of
$\operatorname{Hom}_{S'}(X',Y')$ formed by the homomorphisms compatible with
the canonical descent data on $X',Y'$ (i.e. those whose inverse images under
the two projections from $S''$ to $S'$ are equal). This follows readily from
Theorem 1, Corollary 1, when $Y$ is restricted to be affine over $S$; in the
general case Theorem 1 must be combined with the following result:

\result{LEMMA}{1.2} Let $\alpha:S'\longrightarrow S$ be a faithfully flat and
quasi-compact morphism. Then $S$ identifies with a \emph{quotient topological
space} of $S'$; i.e. every subset $U$ of $S$ such that
$\alpha^{-1}(U)$ is open is itself open.

To complete Theorem 2, one must give effectiveness criteria for a descent datum
on an $S'$-prescheme $X'$ (when $X'$ is not assumed affine over $S'$). First
observe that \emph{such a descent datum need not be effective}, even if $S$ is
the spectrum of a field $k$, $S'$ is the spectrum of a quadratic extension
$k'$ of $k$, and $X'$ is a proper algebraic scheme of dimension 2 over $S'$
(as one can see, following Serre, by using Nagata's nonprojective surface).
\emph{For a descent datum on $X'/S'$ relative to
$\alpha:S'\longrightarrow S$ (faithfully flat and quasi-compact) to be
effective, it is necessary and sufficient that $X'$ be a union of open subsets
$X'_i$, affine over $S'$, that are “stable” under the descent datum on $X'$.}
This certainly holds (whatever $X'/S'$ and its descent datum may be) if the
morphism $\alpha:S'\longrightarrow S$ is \emph{radicial} (i.e. injective, with
residue field extensions that

\src{318}
are radicial). One can also show that it still holds if
$\alpha:S'\longrightarrow S$ is \emph{finite}, and every finite subset of
$X'$ contained in a fiber of $X'$ over $S$ is contained in an open subset of
$X'$ affine over $S$ (this is \emph{Weil's criterion}). In particular, this
holds if $X'/S'$ is quasi-projective; in that case one can show that the
“descended” prescheme $X/S$ is also quasi-projective (and projective if $X'/S'$
is projective). In summary:

\result{THEOREM}{3} Let $\alpha:S'\longrightarrow S$ be a faithfully flat and
quasi-compact morphism of preschemes. If $\alpha$ is \emph{radicial}, it is a
\emph{strict descent morphism}. If $\alpha$ is finite, it is a strict descent
morphism relative to the fibered category of quasi-projective (or projective)
preschemes over preschemes.

\result{REMARKS}{} I do not know whether, in the second assertion above, the
hypothesis that $\alpha$ be a \emph{finite} morphism is really necessary; in
any case one checks “formally” that it may be replaced by the following,
apparently weaker, hypothesis: \emph{for every point of $S$ there exist an open
neighborhood $U$, a $U'$ finite and faithfully flat over $U$, and an
$S$-morphism from $U'$ to $S'$.} A typical case not covered by the preceding
condition is that in which $S=\operatorname{Spec}(A)$,
$S'=\operatorname{Spec}(\widehat A)$, where $A$ is a Noetherian local ring and
$\widehat A$ its completion; another is that in which $S'$ is quasi-finite over
$S$ (i.e. locally isomorphic to an open subset of a finite $S$-scheme) but not
finite. In these two cases, the speaker also does not know the answer to the
following question: if $X$ is an $S$-scheme such that $X'=X\times_SS'$ is
projective over $S'$, is $X$ projective over $S$?\footnote{The 1962 erratum
answers in the negative: a quasi-finite, étale, surjective morphism
$S'\longrightarrow S$, or a morphism
$\operatorname{Spec}(\widehat A)\longrightarrow\operatorname{Spec}(A)$, need
not be a strict descent morphism, even if $A$ is the local ring of an algebraic
curve over an algebraically closed field $k$ and $S=\operatorname{Spec}(A)$.
One may thus find a proper smooth morphism $f:X\longrightarrow S$ making $X$ a
principal bundle over $S$ under an elliptic curve $E$, such that
$f':X'\longrightarrow S'$ is projective but $f$ is not. This is therefore also
an example of a non-isotrivial homogeneous principal bundle under an abelian
scheme.}

\subsection*{2. Application to the descent of certain properties of morphisms}
\addcontentsline{toc}{subsection}{2. Descent of properties of morphisms}

Let $P$ be a class of morphisms of preschemes. Let
$\alpha:S'\longrightarrow S$ be a morphism of preschemes, let
$f:X\longrightarrow Y$ be a morphism of $S$-preschemes, and let
$f':X'\longrightarrow Y'$ be the inverse image of $f$ under $\alpha$. One may
then ask whether “$f'\in P$” implies “$f\in P$.” It appears that the answer is
affirmative in many important cases when $\alpha$ is assumed \emph{faithfully
flat and quasi-compact} (the latter hypothesis sometimes being superfluous).
This follows directly and without difficulty if $P$ is the class of
surjective morphisms, respectively radicial morphisms (these cases following
from the surjectivity of $\alpha$), respectively flat morphisms, faithfully
flat morphisms, smooth morphisms (these cases following from the faithful
flatness of $\alpha$), or morphisms of finite type. Using Theorems 1 and 2 and
Lemma 1.2, one also sees that the same holds when $P$ is one of the following
classes: isomorphisms, open immersions, closed immersions, immersions (if $f$ is
of finite type and $Y$ is locally Noetherian), affine morphisms, finite
morphisms, quasi-finite morphisms, open morphisms, closed morphisms,
homeomorphisms, separated morphisms,

\src{319}
and proper morphisms. The only important unresolved case is that of projective
and quasi-projective morphisms, already noted in the remark of Section 1.

\subsection*{3. Descent by finite faithfully flat morphisms}
\addcontentsline{toc}{subsection}{3. Descent by finite faithfully flat morphisms}

Let $\alpha:S'\longrightarrow S$ be a finite morphism corresponding to a sheaf
of algebras $\mathcal A'$ over $S$ that is locally free of finite type as a
sheaf of modules and is everywhere nonzero. Then $\alpha$ is faithfully flat
and quasi-compact, so the preceding results apply. The datum of a
quasi-coherent sheaf $F'$ over $S'$ is equivalent to that of the quasi-coherent
sheaf $\alpha_*(F')$ over $S$ equipped with its $\mathcal A'$-module structure
(noting that $\mathcal A'=\alpha_*(\mathcal O_{S'})$). For simplicity, this
sheaf over $S$ will also be denoted by $F'$. The two inverse images $p_i^*(F')$
of $F'$ over $S'\times_SS'$ likewise correspond to the quasi-coherent sheaves
of $(\mathcal A'\otimes_{\mathcal O_S}\mathcal A')$-modules
\[
  F'\otimes_{\mathcal O_S}\mathcal A'
  \quad\text{and}\quad
  \mathcal A'\otimes_{\mathcal O_S}F'.
\]
Giving a homomorphism over $S'\times_SS'$ from the first to the second is
equivalent to giving a homomorphism of
$(\mathcal A'\otimes_{\mathcal O_S}\mathcal A')$-modules and, since
$\mathcal A'$ is locally free, this is equivalent to giving a homomorphism of
$(\mathcal A'\otimes\mathcal A')$-modules:
\[
  \mathcal U=
  \operatorname{Hom}_{\mathcal O_S}(\mathcal A',\mathcal A')
  =\mathcal A'\otimes\check{\mathcal A}'
  \longrightarrow
  \operatorname{Hom}_{\mathcal O_S}(F',F'),
\]
i.e. to giving, for every section $\xi$ of $\mathcal U$ over an open subset
$V$, a homomorphism of $\mathcal O_S$-modules
$T_\xi:F'|_V\longrightarrow F'|_V$ satisfying
\begin{equation}
  T_{f\xi}(x)=fT_\xi(x),
  \qquad T_{\xi f}(x)=T_\xi(fx),
  \tag{3.1}
\end{equation}
where $f,x$ are respectively sections of $\mathcal A',F'$ over an open subset
of $S$ contained in $V$. Conditions (i) and (ii) for a descent datum
(A, Section 1(c)) then take the respective forms
\begin{equation}
  T_{1_{\mathcal U}}x=x,
  \qquad\text{i.e.}\qquad
  T_{1_{\mathcal U}}=\operatorname{id}_{F'},
  \tag{3.2}
\end{equation}
\begin{equation}
  T_{\xi\eta}=T_\xi T_\eta.
  \tag{3.3}
\end{equation}
In other words, \emph{a descent datum on $F'$ is equivalent to a representation
of the sheaf of $\mathcal O_S$-algebras
$\mathcal U=\operatorname{Hom}_{\mathcal O_S}(\mathcal A',\mathcal A')$ in the
sheaf of $\mathcal O_S$-algebras
$\operatorname{Hom}_{\mathcal O_S}(F',F')$, satisfying the two linearity
conditions (3.1).} If there is a pairing of quasi-coherent sheaves over $S'$,
\[
  F'_1\times F'_2\longrightarrow F'_3,
\]
(which may be interpreted as a pairing of sheaves of $\mathcal A'$-modules
over $S$),

\src{320}
and gluing data on the $F'_i$, defined by homomorphisms
\[
  T_i:\mathcal U\longrightarrow
  \operatorname{Hom}_{\mathcal O_S}(F'_i,F'_i)
  \qquad (i=1,2,3),
\]
then these data are compatible with the given pairing, in the evident sense,
if and only if the following condition is satisfied.

For every section $\xi$ of $\mathcal U$ over an open subset, write
\[
  \Delta\xi=\sum_i\xi'_i\otimes_{\mathcal A'}\xi''_i
\]
for the section of $\mathcal U\otimes_{\mathcal A'}\mathcal U$ (where
$\mathcal U$ is regarded as an $\mathcal A'$-module through its left module
structure) defined by
\[
  \xi\mathbin{\cdot}(fg)=\sum_i\xi'_i(f)\xi''_i(g)
\]
(where $f$ and $g$ are two sections of $\mathcal A'$ over a smaller open
subset). Then
\begin{equation}
  T^{(3)}_\xi(x\mathbin{\cdot}y)
  =\sum_iT^{(1)}_{\xi'_i}x\mathbin{\cdot}T^{(2)}_{\xi''_i}y
  \tag{3.4}
\end{equation}
for every pair of sections $x,y$ of $F'_1,F'_2$, respectively, over a smaller
open subset.\footnote{The printed text says “of $\mathcal A'$.” The domains of
$T^{(1)}$ and $T^{(2)}$, together with the pairing
$F'_1\times F'_2\longrightarrow F'_3$, require $x$ and $y$ to be sections of
$F'_1$ and $F'_2$, respectively.} (This property may be expressed by saying
that the homomorphisms $T^{(i)}$ are compatible with the diagonal map of
$\mathcal U$ relative to the given pairing.) In particular, formulas
(3.1)--(3.4) allow descent data on a quasi-coherent sheaf of algebras over $S'$
to be interpreted in terms of representations of algebras with diagonal maps,
and hence also (upon restricting to commutative algebras) descent data on an
affine $S'$-scheme.

From this one passes to an analogous interpretation of descent data on an
arbitrary $S'$-prescheme $X'$. Giving such an $X'$ is equivalent to giving a
prescheme $X'$ over $S$ equipped with a homomorphism of $\mathcal O_S$-algebras
\[
  \mathcal A'\longrightarrow\mathcal O_{X'},
\]
and a descent datum on $X'$ is equivalent to giving a homomorphism of sheaves
(compatible with the morphism $h:X'\longrightarrow S'$)
\[
  \mathcal U\longrightarrow
  \operatorname{Hom}_{h^{-1}(\mathcal O_S)}
  (\mathcal O_{X'},\mathcal O_{X'}),
\]
satisfying conditions analogous to (3.1)--(3.4) above.

\textsc{Example 1 (“Weil”).} -- Suppose that $S'/S$ is a Galois étale covering
with Galois group $\Gamma$ (cf. A, Sections 3 and 4(d)). Then a descent datum on
a quasi-coherent sheaf $F'$ over $S'$ (respectively, on an $S'$-prescheme $X'$)
is equivalent to a representation of $\Gamma$ by automorphisms of $(S',F')$
(respectively, of $(S',X')$) compatible with the action of $\Gamma$ on $S'$.
This result

\src{321}
is “formal,” i.e. it can be proved in categorical terms, but from the viewpoint
of this section it also follows from the explicit structure of $\mathcal U$
(equipped with its ring structure, the ring homomorphism
$\mathcal A'\longrightarrow\mathcal U$, and the diagonal map), which is
completely known by virtue of the following result: as a left
$\mathcal A'$-module, $\mathcal U$ admits a basis consisting of those sections
of $\mathcal U$ that correspond to the elements of $\Gamma$.

\textsc{Example 2 (“Cartier”).} -- Let $p$ be a prime number, suppose that
$p\mathcal O_S=0$ (i.e. $\mathcal O_S$ has characteristic $p$), that
$(\mathcal A')^p\subset\mathcal O_S=\mathcal A$ (i.e. $S'/S$ is radicial of
height 1), and that the sheaf of algebras $\mathcal A'$ over $\mathcal A$
locally admits a $p$-basis, i.e. a family $(x_i)$ of sections such that
$\mathcal A'$ is generated as an algebra by the $x_i$, subject only to the
relations $x_i^p=0$. We suppose the set of indices $i$ to be finite, of
cardinality $n$. Let $\mathfrak g$ be the sheaf of $\mathcal A$-derivations of
$\mathcal A'$. It is a locally free sheaf of $\mathcal A'$-modules of rank $n$;
moreover, it is a sheaf of $p$-Lie algebras over $\mathcal A$ (but not over
$\mathcal A'$), satisfying
\begin{equation}
  [X,fY]=X(f)Y+f[X,Y].
  \tag{3.5}
\end{equation}

\textsc{Lemma.} -- $\mathcal U=
\operatorname{Hom}_{\mathcal O_S}(\mathcal A',\mathcal A')$ is generated, as
an $\mathcal O_S$-algebra equipped with an algebra homomorphism
$\mathcal A'\longrightarrow\mathcal U$, by the left $\mathcal A'$-submodule
$\mathfrak g$, subject to the additional relations
\begin{equation}
  \left\{
  \begin{aligned}
    Xf-fX&=X(f),\\
    XY-YX&=[X,Y],\\
    X^p&=X^{(p)}.
  \end{aligned}
  \right.
  \tag{3.6}
\end{equation}

It follows from the preceding lemma that a descent datum on the quasi-coherent
sheaf $F'$ over $S'$ is equivalent to giving, for every $X\in\mathfrak g$, an
$\mathcal O_S$-endomorphism $\overline X$ of $F'$ satisfying
\begin{align}
  \overline{fX}&=f\overline X, \tag{3.7}\\
  \overline X(fx)&=X(f)x+f\overline X(x), \tag{3.8}\\
  \overline{[X,Y]}&=[\overline X,\overline Y], \tag{3.9}\\
  \overline{X^{(p)}}&=\overline X^p. \tag{3.10}
\end{align}
(This is what one might call a linear connection on $F'$, without curvature
and compatible with the $p$-th power.) The notion of a descent datum on an
$S'$-prescheme $X'$ is made explicit in the same way; here relation (3.4) is
replaced by the condition that the $\overline X$ are derivations of
$\mathcal O_{X'}$. Since the morphism $S'\longrightarrow S$ is radicial,
Theorem 3 guarantees that every such descent datum is

\src{322}
effective and hence defines an $S$-prescheme $X$.

We note that no flatness, nonsingularity, or finiteness hypothesis of any kind
has been imposed on $F'$, respectively $X'$.

\subsection*{4. Application to rationality criteria}
\addcontentsline{toc}{subsection}{4. Application to rationality criteria}

Let $X$ be an $S$-prescheme such that the direct image of $\mathcal O_X$ on $S$
is $\mathcal O_S$; this property then remains true under every flat extension
$S'\longrightarrow S$ of the base $S$. If $F$ is an invertible sheaf (i.e.
locally free of rank 1) on $X$, the automorphisms of $F$, which are identified
with the invertible sections of $\mathcal O_X$, correspond bijectively to the
invertible sections of $\mathcal O_S$. Let $s$ now be a section of $X$ over
$S$. We shall refer figuratively to a section of $F$ over $s$, meaning a
section of the invertible sheaf $s^*(F)$ on $S$. It follows from the foregoing
that if $F_i$ ($i=1,2$) are two invertible sheaves on $X$, each equipped with a
section over $s$, and if $F_1$ and $F_2$ are isomorphic, there exists exactly
one isomorphism from $F_1$ to $F_2$ compatible with the sections in question
(i.e. carrying the first to the second). Independently of the section $s$, let
us agree to regard two invertible sheaves $F_1$ and $F_2$ on $X$ as equivalent
if every point of $S$ has an open neighborhood $U$ such that the restrictions
of $F_1$ and $F_2$ to $X|U$ are isomorphic. Then every invertible sheaf $F$ on
$X$ is equivalent to an invertible sheaf $F_1$ equipped with a marked section
over $s$ (take $F_1=F\,s^*(F)^{-1}$), and $F_1$ is determined up to
isomorphism. In other words, the classification of invertible sheaves on $X$
up to equivalence is the same as the classification, up to isomorphism, of
invertible sheaves equipped with a marked section.

Since these properties remain true under a flat extension
$\alpha:S'\longrightarrow S$ of the base (with the section $s$ replaced by its
inverse image $s'$ under $\alpha$), Theorem 1 gives the following conclusion:

\emph{Let the prescheme $X/S$ be as above and admit a section $s$, let
$\alpha:S'\longrightarrow S$ be a faithfully flat and quasi-compact morphism,
and let $F'$ be an invertible sheaf on $X'=X\times_SS'$. Then $F'$ is equivalent
to the inverse image on $X'$ of an invertible sheaf $F$ on $X$ if and only if
its inverse images $p_1^*(F')$ and $p_2^*(F')$ on
$X'\times_XX'=X\times_S(S'\times_SS')$ are equivalent. If so, $F$ is
determined up to equivalence.} (In that case $F'$ is said to be rational over
$S$.)

Applying this principle when $\alpha:S'\longrightarrow S$ is as in Example 1
or Example 2 of the preceding section yields the rationality criteria of Weil
or Cartier. (We note that these authors restrict themselves to the case in
which $S$ and $S'$

\src{323}
are spectra of fields; a fortiori, $S$ is then the spectrum of a local ring,
and the equivalence relation introduced above is simply isomorphism.) In the
first case, $F'$ is rational over $S$ if and only if its transforms under
$\Gamma$ are equivalent to $F'$. To express the rationality criterion in the
second case, consider, in general, the diagonal morphism
\[
  X'\longrightarrow X''=X'\times_XX'
\]
of $X'/X$, the corresponding sheaf of ideals $\mathcal I$ on
$X'\times_XX'$, and the sheaf $\mathcal I/\mathcal I^2$, which is identified
with its inverse image $\Omega^1_{X'/X}$ on $X'$. Since the restrictions to the
diagonal of $F''_i=p_i^*(F')$ ($i=1,2$) are isomorphic (both being isomorphic
to $F'$), i.e. $F''_1(F''_2)^{-1}=F''$ has a trivial restriction to the
diagonal, it follows that the restriction of $F''$ to
$(X'',\mathcal O_{X''}/\mathcal I^2)$ is given, up to isomorphism, by a
well-determined element $\xi$ of
\[
  H^1(X'',\mathcal I/\mathcal I^2)
  =H^1(X',\Omega^1_{X'/X}).
\]
Moreover, in the present case,
\[
  \Omega^1_{X'/X}
  =\Omega^1_{S'/S}\otimes_{\mathcal O_S}\mathcal O_X,
\]
and consequently, if $\Omega^1_{S'/S}$ is locally free over $S$ (as in
Cartier's case), $\xi$ defines a section of
\[
  R^1f'_*(\mathcal O_{X'})\otimes\Omega^1_{S'/S}
\]
over $S'$, called the Atiyah--Cartier class of the invertible sheaf $F'$ on
$X'/S$. Its vanishing is necessary and sufficient for the inverse images of
$F'$ under the two projections from
\[
  (X'',\mathcal O_{X''}/\mathcal I^2)
  =X\times_S(S'',\mathcal O_{S''}/\mathcal J^2)
\]
to $X'$ to be equivalent (where $\mathcal J$ is the sheaf of ideals on
$S''=S'\times_SS'$ defined by the diagonal morphism
$S'\longrightarrow S'\times_SS'$). This vanishing is therefore plainly
necessary for the inverse images of $F'$ on $X''=X\times_SS''$ itself to be
equivalent, and hence also for $F'$ to be equivalent to the inverse image of
an invertible sheaf $F$ on $X$. The Atiyah--Cartier class may also be
interpreted as the obstruction to the existence, locally over $S'$, of a
connection on $F'$ relative to the derivations of $X'/X$; such a connection is
further determined once one knows the derivations of $F'$ corresponding to the
natural prolongations to $X'$ of the derivations of $S'/S$. From this and the
developments in the preceding section, one readily concludes that in the case
of Example 2 of that section, and when $X/S$ admits a section, the vanishing of
the Atiyah--Cartier class is also sufficient for $F'$ to be rational over $S$.

\subsection*{5. Application to restricting the base in an abelian scheme}
\addcontentsline{toc}{subsection}{5. Restricting the base in an abelian scheme}

Let $S$ be a prescheme. An abelian scheme over $S$ is a scheme $X$ smooth and
proper

\src{324}
over $S$ whose fibers at the points $x\in S$ are schemes of abelian varieties
over the fields $\kappa(x)$. Suppose that $S$ is Noetherian and regular (i.e.
its local rings are regular). Using Murre's connectedness theorem [4], one can
then show (at least in the equal-characteristic case, in which the cited
theorem is presently proved) that every rational section of $X$ over $S$ is
defined everywhere (i.e. is a section), generalizing a classical theorem of
Weil. More generally, it follows that if $X'$ is smooth over $S$, then every
rational $S$-map from $X'$ to $X$ is defined everywhere. This gives the
following result, which generalizes a result of Chow--Iang:\footnote{The
printed text says “Igusa--Lang”; the 1962 erratum prescribes “Chow--Iang.”}
Let $S$ be Noetherian and regular, and let $K$ denote its ring of rational
functions (a direct product of fields). Let $X$ be an abelian scheme over $K$.
If $X$ is isomorphic to a $K$-scheme of the form
$X_0\times_S\operatorname{Spec}(K)$, where $X_0$ is an abelian scheme over
$S$, then $X_0$ is determined up to a unique isomorphism.

Using the preceding uniqueness result, one sees that the problem of
restricting the base in $X$ is local on $S$ (and hence that it suffices to know
how to carry out the restriction over the $\operatorname{Spec}(\mathcal O_x)$,
with $x\in S$). In the same way, if $\alpha:S'\longrightarrow S$ is a smooth
morphism of finite type, if $K'$ is the ring of rational functions of $S'$,
and if $X\otimes_KK'$ is of the form
$X'_0\times_{S'}\operatorname{Spec}(K')$, then $X'_0$ carries a canonical
descent datum relative to $\alpha$. In view of Theorem 3, it follows that:

\result{PROPOSITION}{5.1} Let $S$ be an irreducible Noetherian regular
prescheme with rational function field $K$, let $K'$ be a finite extension of
$K$ unramified over $S$, let $S'$ be the normalization of $S$ in $K'$ (thus an
étale covering of $S$), and let $X$ be an abelian scheme over $K$ such that
$X\otimes_KK'$ is of the form
$X'_0\times_{S'}\operatorname{Spec}(K')$, where $X'_0$ is a projective abelian
scheme over $S'$. Then $X$ is of the form
$X_0\times_S\operatorname{Spec}(K)$, where $X_0$ is a projective abelian scheme
over $S$.

\result{REMARKS}{} The speaker does not know whether the hypothesis that
$S'\longrightarrow S$ is a surjective étale covering (which permits the use of
Theorem 3) can be replaced by the hypothesis that it is a surjective smooth
morphism of finite type (even if it is assumed to be étale), or whether the
proposition remains valid without assuming $X'_0$ projective over $S'$ (a
condition that may perhaps hold automatically).

\subsection*{6. Application to criteria for local triviality and isotriviality}
\addcontentsline{toc}{subsection}{6. Criteria for local triviality and isotriviality}

Let $S$ be a prescheme, let $G$ be a “prescheme in groups” over $S$, and let
$P$ be a prescheme over $S$ on which “$G$ acts” (on the right). We say that
$P$ is formally principal homogeneous under $G$ if the familiar morphism
\[
  G\times_SP\longrightarrow P\times_SP
\]

\src{325}
induced by the action of $G$ on $P$ is an isomorphism. Henceforth we assume
that $G$ is flat over $S$ (and therefore faithfully flat over $S$), and reserve
the name principal homogeneous bundle under $G$ for a formally principal
homogeneous bundle $P$ that is faithfully flat and quasi-compact over $S$.
Equivalently, there is a faithfully flat and quasi-compact extension
$S'\longrightarrow S$ of the base $S$ such that the formally principal
homogeneous bundle
\[
  P'=P\times_SS'
\]
under $G'=G\times_SS'$ is trivial, i.e. isomorphic to $G'$ (i.e. admits a
section); one may in particular take $S'=P$. Note also that if $S$ is locally
Noetherian, then the hypothesis that $P$ be faithfully flat is equivalent to
the hypothesis that
\[
  \widehat P_s=P\times_S\operatorname{Spec}(\widehat{\mathcal O}_s)
\]
be faithfully flat over $\widehat{\mathcal O}_s$ for every $s\in S$ (where
$\widehat{\mathcal O}_s$ denotes the completion of the local ring
$\mathcal O_s$), since $\widehat{\mathcal O}_s$ is faithfully flat over
$\mathcal O_s$. Moreover, if $P$ is of finite type over the locally Noetherian
prescheme $S$, the set of points $s$ satisfying the preceding condition is
constructible. Thus, if $S$ is a “Jacobson prescheme” (for example, a scheme
of finite type over a field, or more generally the spectrum of a Jacobson
ring), it suffices to check the condition at the closed points of $S$. This
reduces us to the case
in which the base is the spectrum of a complete local ring $A$. When
$S=\operatorname{Spec}(A)$, where $A$ is a complete Noetherian local ring, and
$P$ is of finite type over $S$, the faithful flatness of $P/S$ is also
equivalent to the existence of an $S'$ finite and flat over $S$ such that $P'$
is trivial; if in addition $G$ is smooth over $S$, one may take $S'$ étale
over $S$. Consequently, if the residue field of $A$ is moreover algebraically
closed (the “geometric case”), $P$ is faithfully flat over $A$ if and only if
it is trivial. Thus, if $S$ is an algebraic prescheme over an algebraically
closed field and $G$ is smooth and of finite type over $S$, the faithful
flatness condition over $S$ is equivalent to Serre's analytic triviality
condition (SLF) ([6], pp. 1--12).

One may introduce other, stronger types of conditions on $P$ having the nature
of a “local triviality.” We say in particular that $P$ is isotrivial
(respectively, strictly isotrivial) if, for every $s\in S$, there exist an open
neighborhood $U$ of $s$ and a finite faithfully flat morphism (respectively, a
surjective étale covering) $U'\longrightarrow U$ such that
$P'=P\times_SU'$ is trivial. (Our terminology differs from Serre's [6], who
calls locally isotrivial what we call strictly isotrivial.) Strict
isotriviality is especially useful when $G$ is smooth over $S$, but is an
inadequate notion in the other cases.

If $G$ is affine over $S$, every principal homogeneous bundle $P$ under $G$ is
affine by Section 2, and Theorem 2 therefore makes it possible to “descend”

\src{326}
such bundles along faithfully flat and quasi-compact morphisms. In particular,
take $G=\operatorname{Gl}(n)_S$, defined by the requirement that the functor on
$S$-preschemes
\[
  S'\longmapsto\operatorname{Hom}_S(S',G)
\]
(with values in the category of groups) be identified with the functor
\[
  \operatorname{Gl}(n)(S')
  =\operatorname{Gl}\bigl(n,H^0(S',\mathcal O_{S'})\bigr)
\]
of A, Section 4(e). Use the facts (i) that every principal homogeneous bundle
under $G$ (respectively, every locally free sheaf of rank $n$ over $S$) becomes
isomorphic to the “trivial” object $G$ (respectively, $\mathcal O_S^n$) after a
suitable faithfully flat and quasi-compact extension of $S$; (ii) that objects
of the type in question (principal homogeneous bundles under $G$,
respectively, locally free sheaves of rank $n$) can be descended along such
morphisms; and (iii) that the automorphism group of the trivial bundle over an
$S'/S$ is functorially isomorphic to the automorphism group of the trivial
locally free sheaf of rank $n$ over $S'$. One concludes “formally” that giving
a principal homogeneous bundle with group $G$ over $S$ (or over an $S'/S$) is
“the same thing” as giving a locally free sheaf of rank $n$ there. (More
precisely, there is an equivalence of fibered categories.) In particular:

\result{PROPOSITION}{6.1} Every principal homogeneous bundle with group
$\operatorname{Gl}(n)_S$ is locally trivial.

Known arguments give the same result for structure groups such as
$\operatorname{Sl}(n)_S$, $\operatorname{Sp}(n)_S$, and products of such
groups. It also follows that if $F$ is a closed subgroup of
$G=\operatorname{Gl}(n)_S$, flat over $S$, such that the quotient $G/F$ exists,
and if $G$ is an isotrivial (respectively, strictly isotrivial) principal
homogeneous bundle over $G/F$ with group $F\times_S(G/F)$, then every principal
homogeneous bundle with group $F$ is isotrivial (respectively, strictly
isotrivial). This applies to all the “linear groups” over $S$ used so far, and
in particular when $G=S\times_k\Gamma$, where $S$ is a prescheme over the field
$k$ and $\Gamma$ is a linear group in the classical sense, and in particular
a smooth group, over $k$. For such groups this therefore resolves a question
of Serre (loc. cit.).

We also note that for most known smooth groups over $S$, whether linear or not,
and in any case for those of the form $S\times_k\Gamma$ above, one can show
that every isotrivial principal homogeneous bundle is strictly isotrivial.
This resolves in particular another question of Serre (loc. cit., pp. 1--14),
since a principal homogeneous bundle obtained by descent à la Cartier (cf.
Section 3, Example 2) is manifestly isotrivial.

\result{REMARK}{} One of the essential difficulties in these questions (apart
from the question of the existence of the quotient schemes $G/F$) is the lack
of effectivity criteria for a descent datum along a faithfully flat nonfinite
morphism.

\src{327}
\begin{center}
  \textbf{BIBLIOGRAPHY}
\end{center}

\begingroup
\small
\begin{description}
  \item[{[1]}] Dieudonné (J.) and Grothendieck (Alexander). --- \emph{Éléments
    de géométrie algébrique}, forthcoming in the Publications mathématiques de
    l'Institut des Hautes Études Scientifiques.
  \item[{[2]}] Grauert (H.) und Remmert (R.). --- \emph{Komplexe Räume}, Math.
    Annalen, vol. 136, 1958, pp. 245--318.
  \item[{[3]}] Grothendieck (Alexander). --- \emph{Géométrie formelle et
    géométrie algébrique}. Séminaire Bourbaki, vol. 11, 1958/59, no. 182.
  \item[{[4]}] Murre (J. P.). --- \emph{On a connectedness theorem for a
    birational transformation at a simple point}, Amer. J. Math., vol. 80,
    1958, pp. 3--15.
  \item[{[5]}] Serre (Jean-Pierre). --- \emph{Géométrie algébrique et géométrie
    analytique}, Ann. Inst. Fourier Grenoble, vol. 6, 1955--56, pp. 1--42.
  \item[{[6]}] Serre (Jean-Pierre). --- \emph{Espaces fibrés algébriques},
    Séminaire Chevalley, vol. 2, 1958: Anneaux de Chow et applications, no. 1.
\end{description}
\endgroup
